\chapter{Matrizes e grafos}\label{ch:g12:matrix}

Uma \omterm{def:g12:matrix:matrix}{matriz} é uma tabela retangular de números, somada e multiplicada por
regras concebidas para que a álgebra das \omterm{def:g12:matrix:matrix}{matrizes} represente a composição
de transformações lineares. As \omterm{def:g12:matrix:matrix}{matrizes} resolvem \omterm{def:g10:lines:system}{sistemas lineares},
governam \omterm{def:g12:seq:sequence}{sequências} recorrentes acopladas e contam passeios em redes ---
a matemática por trás dos mecanismos de busca e dos algoritmos de caminho
\omterm{def:g10:functions:extrema}{mínimo}.

\section{Álgebra das matrizes}

\begin{definition}[Matriz]\label{def:g12:matrix:matrix}
Uma \emph{matriz $m \times n$}\index{matriz} é uma tabela de \omterm{def:g10:numbers:sets}{números reais}
com $m$ linhas e $n$ colunas:
$A = (a_{ij})$, em que $a_{ij}$ é a entrada da linha $i$, coluna $j$.
Duas matrizes de mesmo tamanho são somadas entrada a entrada, e
$\lambda A = (\lambda a_{ij})$.
\end{definition}

\begin{definition}[Produto de matrizes]\label{def:g12:matrix:product}
Sejam $A$ de tamanho $m \times n$ e $B$ de tamanho $n \times p$. O produto
$AB$ é a \omterm{def:g12:matrix:matrix}{matriz} $m \times p$ cuja entrada $(i,j)$ é
\[
(AB)_{ij} = \sum_{k=1}^{n} a_{ik} b_{kj}
\]
(a regra ``linha $i$ de $A$ vezes coluna $j$ de $B$'').
\end{definition}

\begin{example}
$\begin{pmatrix} 1 & 2\\ 3 & 4\end{pmatrix}
\begin{pmatrix} 0 & 1\\ 1 & 1\end{pmatrix}
= \begin{pmatrix} 2 & 3\\ 4 & 7\end{pmatrix}$,
\quad ao passo que \quad
$\begin{pmatrix} 0 & 1\\ 1 & 1\end{pmatrix}
\begin{pmatrix} 1 & 2\\ 3 & 4\end{pmatrix}
= \begin{pmatrix} 3 & 4\\ 4 & 6\end{pmatrix}$:
\emph{a multiplicação de \omterm{def:g12:matrix:matrix}{matrizes} não é comutativa}.
\end{example}

\begin{proposition}[Regras da álgebra das matrizes]\label{prop:g12:matrix:rules}
Sempre que os tamanhos tornem os produtos legítimos:
\[
(AB)C = A(BC), \qquad A(B + C) = AB + AC, \qquad (A+B)C = AC + BC,
\]
e a \emph{\omterm{def:g12:matrix:matrix}{matriz} identidade} $I_n$ (uns na diagonal, zeros no resto)
satisfaz $I_m A = A I_n = A$ para $A$ de tamanho $m \times n$.
\end{proposition}

\begin{proof}
Todas são verificações entrada a entrada a partir da
\cref{def:g12:matrix:product}; a associatividade, a única não trivial,
resume-se a trocar duas somas finitas:
$\bigl((AB)C\bigr)_{ij} = \sum_l \left(\sum_k a_{ik}b_{kl}\right) c_{lj}
= \sum_k a_{ik} \left(\sum_l b_{kl} c_{lj}\right)
= \bigl(A(BC)\bigr)_{ij}$.
\end{proof}

\begin{definition}[Inversa]\label{def:g12:matrix:inverse}
Uma \omterm{def:g12:matrix:matrix}{matriz} quadrada $A$ de tamanho $n$ é \emph{invertível}\index{matriz!inversa}
se existe uma \omterm{def:g12:matrix:matrix}{matriz} $B$ com $AB = BA = I_n$; esse $B$ é então único e se
escreve $A^{-1}$.
\end{definition}

\begin{proposition}[Inversa de uma matriz $2\times2$]\label{prop:g12:matrix:inverse2x2}
Sejam $A = \begin{pmatrix} a & b\\ c & d\end{pmatrix}$ e
$\det A = ad - bc$ (o \emph{determinante}\index{determinante}). Então $A$ é
\omterm{def:g12:matrix:inverse}{invertível} se, e somente se, $\det A \neq 0$, e nesse caso
\[
A^{-1} = \frac{1}{ad - bc}\begin{pmatrix} d & -b\\ -c & a\end{pmatrix}.
\]
\end{proposition}

\begin{proof}
Um cálculo dá
$A \begin{pmatrix} d & -b\\ -c & a\end{pmatrix}
= \begin{pmatrix} d & -b\\ -c & a\end{pmatrix} A = (ad - bc) I_2$; se
$ad - bc \neq 0$, divida. Reciprocamente, se $ad - bc = 0$, as colunas de
$A$ são proporcionais, e o mesmo vale para as colunas de $AB$ qualquer que
seja $B$; ora, as colunas de $I_2$ não são proporcionais, de modo que
nenhum $B$ pode satisfazer $AB = I_2$.
\end{proof}

\begin{method}[Sistemas lineares]\label{met:g12:matrix:systems}
O sistema
$\begin{cases} ax + by = e\\ cx + dy = f \end{cases}$
é a \omterm{def:g10:algebra:equation}{equação} matricial $AX = Y$ com
$X = \begin{pmatrix} x \\ y\end{pmatrix}$,
$Y = \begin{pmatrix} e \\ f\end{pmatrix}$. Se $\det A \neq 0$, sua única
solução é $X = A^{-1}Y$. O mesmo formalismo trata $n$ \omterm{def:g10:algebra:equation}{equações} a $n$
incógnitas.
\end{method}

\section{Potências de matrizes e sequências recorrentes}

\begin{definition}\label{def:g12:matrix:power}
Para uma \omterm{def:g12:matrix:matrix}{matriz} quadrada $A$ e $k \in \N$, $A^k = A \times \dots \times A$
($k$ fatores), com $A^0 = I$.
\end{definition}

\begin{method}[Casos diagonal-mais-nilpotente e diagonalizável]\label{met:g12:matrix:powers}
Duas maneiras usuais de calcular $A^k$:
\begin{itemize}
  \item Se $A = \lambda I + N$ com $N^2 = 0$, o teorema binomial
        (válido aqui porque $I$ e $N$ comutam) se reduz a dois termos:
        $A^k = \lambda^k I + k \lambda^{k-1} N$.
  \item Se encontramos $P$ \omterm{def:g12:matrix:inverse}{invertível} com $A = PDP^{-1}$ e $D$ diagonal,
        então $A^k = P D^k P^{-1}$, e $D^k$ se calcula entrada a entrada.
        (Encontrar tal $P$ de maneira sistemática é a teoria da
        \emph{diagonalização}, desenvolvida na graduação; neste nível $P$ é
        dado.)
\end{itemize}
\end{method}

\begin{example}[Sequências acopladas]\label{ex:g12:matrix:coupled}
Sejam $u_{n+1} = 3u_n + v_n$ e $v_{n+1} = u_n + 3v_n$. Pondo
$X_n = \begin{pmatrix} u_n\\ v_n \end{pmatrix}$ e
$A = \begin{pmatrix} 3 & 1\\ 1 & 3\end{pmatrix}$, obtemos
$X_{n+1} = AX_n$, logo $X_n = A^n X_0$. As \omterm{def:g12:seq:sequence}{sequências} auxiliares
$s_n = u_n + v_n$ e $d_n = u_n - v_n$ satisfazem $s_{n+1} = 4s_n$ e
$d_{n+1} = 2d_n$, de modo que $s_n = 4^n s_0$, $d_n = 2^n d_0$ e
\[
u_n = \frac{4^n(u_0+v_0) + 2^n(u_0-v_0)}{2}, \qquad
v_n = \frac{4^n(u_0+v_0) - 2^n(u_0-v_0)}{2}.
\]
(Nos bastidores: $(1,1)$ e $(1,-1)$ são direções de autovetores de $A$.)
\end{example}

\section{Grafos e passeios}

\begin{definition}[Grafo, matriz de adjacência]\label{def:g12:matrix:graph}
Um \emph{grafo}\index{grafo} é formado por vértices $1, 2, \dots, n$ e
arestas ligando certos pares de vértices (pares ordenados, no caso de um
grafo \emph{orientado}). Sua \emph{matriz de adjacência}\index{matriz de
adjacência} é a \omterm{def:g12:matrix:matrix}{matriz} $M$ de tamanho $n \times n$ com $m_{ij} = 1$ se há
uma aresta de $i$ para $j$, e $0$ caso contrário. Um \emph{passeio} de
comprimento $k$ de $i$ a $j$ é uma \omterm{def:g12:seq:sequence}{sequência} de $k$ arestas consecutivas
que leva de $i$ a $j$.
\end{definition}

\begin{omfigure}
\begin{tikzpicture}[baseline=(current bounding box.center)]
  \node[circle, draw, inner sep=2.5pt] (v1) at (0,0) {$1$};
  \node[circle, draw, inner sep=2.5pt] (v2) at (2.6,1.2) {$2$};
  \node[circle, draw, inner sep=2.5pt] (v3) at (2.6,-1.2) {$3$};
  \draw[-Stealth, omDef, thick] (v1) -- (v2);
  \draw[-Stealth, omDef, thick] (v2) -- (v3);
  \draw[-Stealth, omDef, thick] (v1) to[bend right=18] (v3);
  \draw[-Stealth, omDef, thick] (v3) to[bend right=18] (v1);
\end{tikzpicture}
\qquad\qquad
$M = \begin{pmatrix} 0 & 1 & 1\\ 0 & 0 & 1\\ 1 & 0 & 0 \end{pmatrix}$

{\small Um \omterm{def:g12:matrix:graph}{grafo} orientado e sua \omterm{def:g12:matrix:graph}{matriz de adjacência}
(\cref{exo:g12:matrix:6}): $m_{ij} = 1$ exatamente quando há uma aresta
de $i$ para $j$.}
\end{omfigure}

\begin{theorem}[Contagem de passeios]\label{thm:g12:matrix:walks}
O número de passeios de comprimento $k$ do vértice $i$ ao vértice $j$ é a
entrada $(i,j)$ de $M^k$.
\end{theorem}

\begin{proof}
Indução em $k$. Para $k = 1$, é a definição de $M$. Suponha a afirmação
verdadeira para $k$. Um passeio de comprimento $k+1$ de $i$ a $j$ é um
passeio de comprimento $k$ de $i$ até algum vértice $l$, seguido de uma
aresta de $l$ a $j$; pelos princípios aditivo e multiplicativo, seu número
é
\[
\sum_{l=1}^{n} \bigl(M^k\bigr)_{il}\, m_{lj} = \bigl(M^{k+1}\bigr)_{ij}.
\qedhere
\]
\end{proof}

\begin{example}
Para o \omterm{def:g12:matrix:graph}{grafo} triangular ($3$ vértices, todos os pares ligados),
$M = \begin{pmatrix} 0&1&1\\ 1&0&1\\ 1&1&0\end{pmatrix}$ e
$M^2 = \begin{pmatrix} 2&1&1\\ 1&2&1\\ 1&1&2\end{pmatrix}$: de cada
vértice há $2$ passeios de comprimento $2$ de volta a ele mesmo (por um ou
outro vizinho) e $1$ até cada um dos outros vértices.
\end{example}

\section{Exercícios}

\begin{exercise}[$\star$]\label{exo:g12:matrix:1}
Sejam $A = \begin{pmatrix} 1 & 2\\ 0 & 1 \end{pmatrix}$ e
$B = \begin{pmatrix} 2 & 0\\ 1 & 1 \end{pmatrix}$. Calcule $A + B$, $AB$,
$BA$ e $A^2$.
\end{exercise}

\begin{exercise}[$\star$]\label{exo:g12:matrix:2}
Determine se as \omterm{def:g12:matrix:matrix}{matrizes} seguintes são \omterm{def:g12:matrix:inverse}{invertíveis} e calcule as inversas
quando existirem:
\[
A = \begin{pmatrix} 2 & 5\\ 1 & 3\end{pmatrix}, \qquad
B = \begin{pmatrix} 3 & 6\\ 2 & 4\end{pmatrix}.
\]
\end{exercise}

\begin{exercise}[$\star$]\label{exo:g12:matrix:3}
Resolva por inversão matricial o sistema
$\begin{cases} 2x + 5y = 1\\ x + 3y = 2 . \end{cases}$
\end{exercise}

\begin{exercise}[$\star\star$]\label{exo:g12:matrix:4}
Seja $A = \begin{pmatrix} 2 & 1\\ 0 & 2\end{pmatrix} = 2I + N$ com
$N = \begin{pmatrix} 0 & 1\\ 0 & 0 \end{pmatrix}$.
\begin{enumerate}
  \item Verifique que $N^2 = 0$ e que $I$ e $N$ comutam.
  \item Deduza $A^k$ para todo $k \in \N$ e confira a fórmula para $k=2$
        por cálculo direto.
\end{enumerate}
\end{exercise}

\begin{exercise}[$\star\star$]\label{exo:g12:matrix:5}
Sejam $A = \begin{pmatrix} 0 & 1\\ 1 & 1\end{pmatrix}$ e $F_n$ a \omterm{def:g12:seq:sequence}{sequência}
de Fibonacci ($F_0 = 0$, $F_1 = 1$, $F_{n+2} = F_{n+1} + F_n$). Mostre por
indução que, para $n \geq 1$,
\[
A^n = \begin{pmatrix} F_{n-1} & F_n\\ F_n & F_{n+1}\end{pmatrix},
\]
e deduza a identidade $F_{n+1}F_{n-1} - F_n^2 = (-1)^n$.
(Dica: os \omterm{def:g11:vect:det}{determinantes} se multiplicam: $\det(MN) = \det M \det N$, o que
você pode verificar para \omterm{def:g12:matrix:matrix}{matrizes} $2\times2$.)
\end{exercise}

\begin{exercise}[$\star\star$]\label{exo:g12:matrix:6}
Um \omterm{def:g12:matrix:graph}{grafo} orientado sobre os vértices $\{1, 2, 3\}$ tem as arestas
$1\to2$, $2\to3$, $3\to1$ e $1\to3$.
\begin{enumerate}
  \item Escreva a \omterm{def:g12:matrix:graph}{matriz de adjacência} $M$ e calcule $M^2$ e $M^3$.
  \item Quantos passeios de comprimento $3$ vão de $1$ a $1$? Liste-os.
\end{enumerate}
\end{exercise}

\begin{exercise}[$\star\star$]\label{exo:g12:matrix:7}
Uma empresa de compartilhamento de carros movimenta veículos entre duas
cidades $A$ e $B$. A cada semana, $80\%$ dos carros em $A$ permanecem em
$A$ e $20\%$ vão para $B$; $30\%$ dos carros em $B$ vão para $A$ e $70\%$
permanecem. Sejam $a_n, b_n$ as proporções da frota em cada cidade.
\begin{omfigure}
\begin{tikzpicture}
  \node[circle, draw, inner sep=3pt] (A) at (0,0) {$A$};
  \node[circle, draw, inner sep=3pt] (B) at (3.2,0) {$B$};
  \draw[-Stealth, omDef, thick] (A) to[bend left=25]
    node[above, font=\small] {$0.2$} (B);
  \draw[-Stealth, omDef, thick] (B) to[bend left=25]
    node[below, font=\small] {$0.3$} (A);
  \draw[-Stealth, omDef, thick] (A) to[loop left]
    node[left, font=\small] {$0.8$} (A);
  \draw[-Stealth, omDef, thick] (B) to[loop right]
    node[right, font=\small] {$0.7$} (B);
\end{tikzpicture}
\end{omfigure}
\begin{enumerate}
  \item Escreva $X_{n+1} = MX_n$ com
        $X_n = \begin{pmatrix} a_n\\ b_n\end{pmatrix}$ e identifique $M$.
  \item Determine as proporções de equilíbrio (resolva $MX = X$ com
        $a + b = 1$).
  \item Mostre que $c_n = a_n - 0.6$ satisfaz $c_{n+1} = 0.5\,c_n$ e
        conclua que a \omterm{def:g11:prob:rv}{distribuição} da frota \omterm{def:g12:seq:limit}{converge} para o equilíbrio.
\end{enumerate}
\end{exercise}

\begin{exercise}[$\star\star\star$]\label{exo:g12:matrix:8}
Sejam $A = \begin{pmatrix} 3 & 1\\ 1 & 3 \end{pmatrix}$,
$P = \begin{pmatrix} 1 & 1\\ 1 & -1 \end{pmatrix}$.
\begin{enumerate}
  \item Calcule $P^{-1}$, depois $D = P^{-1}AP$, e verifique que $D$ é
        diagonal.
  \item Deduza uma fórmula fechada para $A^n$ e compare com o
        \cref{ex:g12:matrix:coupled}.
\end{enumerate}
\end{exercise}

\section{Problema: a matriz que conhece Fibonacci (e o tempo)}

\begin{problem}[{Problema de fim de semana --- uma única matriz
$2 \times 2$ carrega toda a sequência de Fibonacci, uma matriz de
Markov prevê o tempo a longo prazo e um autovetor vale um bilhão
de dólares}]
\label{pb:g12:matrix:1}
Uma \omterm{def:g12:matrix:matrix}{matriz} é uma máquina que come um estado e devolve o
seguinte --- e suas \emph{potências} guardam, portanto, futuros
inteiros. Este problema abre com a espantosa \omterm{def:g12:matrix:matrix}{matriz} cujas
potências listam os números de Fibonacci (e demonstram suas
identidades em uma linha cada), depois faz o tempo rodar como
uma cadeia de Markov até seu regime estacionário e fecha com o
autovetor sobre o qual se construiu um mecanismo de busca
(\cref{thm:g12:matrix:walks},
\cref{met:g12:matrix:powers}).

\medskip
\noindent\textbf{Parte I --- Fluência.}
\begin{enumerate}
  \item Com $A = \begin{pmatrix} 1 & 2\\ 3 & 4\end{pmatrix}$
        e $B = \begin{pmatrix} 0 & 1\\ 1 & 0\end{pmatrix}$:
        calcule $AB$ e $BA$. Veredicto sobre a comutatividade?
  \item Inverta $\begin{pmatrix} 2 & 1\\ 5 & 3\end{pmatrix}$
        (\cref{prop:g12:matrix:inverse2x2}) e use a inversa
        para resolver $2x + y = 4$, $5x + 3y = 7$.
  \item Seja $N = \begin{pmatrix} 0 & 1\\ 0 & 0\end{pmatrix}$:
        calcule $N^2$ e deduza que
        $(I + N)^n = I + nN$ para todo $n$.
  \item O \omterm{def:g12:matrix:graph}{grafo} triangular (três vértices, todos os pares
        ligados): escreva sua \omterm{def:g12:matrix:graph}{matriz de adjacência} $A$, calcule
        $A^3$ e interprete as entradas da diagonal
        (\cref{thm:g12:matrix:walks}).
  \item Para $D = \begin{pmatrix} 2 & 0\\ 0 & \frac12
        \end{pmatrix}$: dê $D^n$ e seu comportamento quando
        $n \to \infty$.
\end{enumerate}

\medskip
\noindent\textbf{Parte II --- A \omterm{def:g12:matrix:matrix}{matriz} de Fibonacci.}
Seja $F = \begin{pmatrix} 1 & 1\\ 1 & 0\end{pmatrix}$ e sejam
$F_1 = F_2 = 1, F_3 = 2, \dots$ os números de Fibonacci do
\cref{pb:g11:seq:1}.
\begin{enumerate}[resume]
  \item Calcule $F^2$, $F^3$, $F^4$ e conjecture a forma geral
        de $F^n$ em termos dos números de Fibonacci.
  \item Demonstre a conjectura
        $F^n = \begin{pmatrix} F_{n+1} & F_n\\ F_n & F_{n-1}
        \end{pmatrix}$ por indução.
  \item Tome os \omterm{def:g11:vect:det}{determinantes} dos dois lados (o \omterm{def:g11:vect:det}{determinante} de
        um produto é o produto dos \omterm{def:g11:vect:det}{determinantes} --- verifique
        para \omterm{def:g12:matrix:matrix}{matrizes} $2 \times 2$ se nunca viu isso): deduza a
        \emph{identidade de Cassini}
        $F_{n+1}F_{n-1} - F_n^2 = (-1)^n$ --- o motor do
        quadrado que some, demonstrado em uma linha.
  \item De $F^{m+n} = F^m F^n$, leia as entradas superiores
        direitas e obtenha a \emph{fórmula de adição}
        \[
        F_{m+n} = F_{m+1} F_n + F_m F_{n-1} .
        \]
        Verifique-a para $m = n = 3$.
  \item Deduza da fórmula de adição (indução em $k$)
        que $F_n$ \omterm{def:g12:arith:divides}{divide} $F_{kn}$, e confira em
        $F_3 \mid F_6$ e $F_3 \mid F_9$.
  \item Para calcular $F_{100}$, não é preciso multiplicar $100$
        \omterm{def:g12:matrix:matrix}{matrizes}: eleve ao quadrado repetidamente ($F^2, F^4,
        F^8, \dots$) e combine. Quantas multiplicações de
        \omterm{def:g12:matrix:matrix}{matrizes} bastam, e que antigo truque de multiplicação do
        volume do ensino fundamental é este, promovido a
        \omterm{def:g12:matrix:matrix}{matrizes}?
\end{enumerate}

\medskip
\noindent\textbf{Parte III --- A máquina do tempo.}
Em certa cidade: depois de um dia de sol, o seguinte é de sol
com \omterm{def:g10:proba:distribution}{probabilidade} $0.8$; depois de um dia de chuva, é de sol com
\omterm{def:g10:proba:distribution}{probabilidade} $0.4$. Codifique a \omterm{def:g11:prob:rv}{distribuição} do dia como uma
coluna $\binom{p_{\text{sol}}}{p_{\text{chuva}}}$ e a evolução
por
\[
M = \begin{pmatrix} 0.8 & 0.4\\ 0.2 & 0.6 \end{pmatrix}.
\]
\begin{enumerate}[resume]
  \item Verifique que cada coluna de $M$ soma $1$ e diga por que
        toda máquina do tempo tem de ter essa propriedade.
  \item Hoje faz sol. Calcule a previsão para amanhã e para
        depois de amanhã.
  \item Encontre o \emph{regime estacionário}: a \omterm{def:g11:prob:rv}{distribuição}
        $v$ com $Mv = v$ (e entradas somando $1$). Que fração
        dos dias é de sol a longo prazo?
  \item Parta de um dia de chuva, $\binom01$, e aplique $M$
        quatro vezes, acompanhando a distância ao regime
        estacionário em cada passo. Por que fator a diferença
        encolhe a cada passo --- e que tipo de convergência é
        esta?
  \item PageRank em miniatura: três páginas, com os links
        $A \to B$, $A \to C$, $B \to C$, $C \to A$. Um
        navegante aleatório segue um link de saída escolhido ao
        acaso, uniformemente. Escreva a \omterm{def:g12:matrix:matrix}{matriz} de transição,
        encontre o regime estacionário e ordene as páginas.
  \item Interprete a ordenação: por que $C$ pontua tão alto
        quanto $A$ apesar de receber links de menos páginas ---
        o que o regime estacionário realmente mede? (O PageRank
        de verdade acrescenta um fator de amortecimento para
        becos sem saída e saltos; a ideia do autovetor é
        exatamente esta.)
\end{enumerate}

\medskip
\noindent\textbf{Parte IV --- Dividendos diagonais.}
\begin{enumerate}[resume]
  \item Duas grandezas acopladas obedecem a
        $u_{n+1} = 3u_n + v_n$, $v_{n+1} = u_n + 3v_n$, ou
        seja, à \omterm{def:g12:matrix:matrix}{matriz} $A$ do \cref{exo:g12:matrix:8}. Usando a
        diagonalização daquele exercício
        ($D = \operatorname{diag}(4, 2)$), dê a fórmula fechada
        para $u_n$ quando $u_0 = 1$, $v_0 = 0$, e confira-a
        contra o cálculo direto para $n = 1, 2, 3$.
  \item Em uma ou duas frases: o que a diagonalização
        \emph{faz} com um sistema acoplado --- e em que sentido
        o regime estacionário de Markov da questão 14 é
        também uma história de autovetores?
  \item Finale --- as três faces da \omterm{def:g12:matrix:matrix}{matriz} neste fim de semana:
        contabilidade (sistemas e inversas), combinatória
        (passeios e links contados por potências) e evolução
        (Fibonacci, o tempo, a web --- futuros lidos em
        direções próprias). Uma frase para cada, mais o
        ponteiro adiante: a álgebra linear dos volumes de
        graduação transforma cada uma dessas faces em uma
        teoria.
\end{enumerate}
\end{problem}
